\section{讨论}

\subsection{预测问题的界定与材料信息的实际增量}

% Word body block 99
本研究把历史模型评价和新增实验观测置于同一研究问题下，同时保留两者不同的证据任务。历史分组结果显示，当前基础输入对跨文献存活率的预测仍有较大误差，更多化学字段、显式剂量—时间交互与固定预算迁移没有稳定降低误差；新增实验则提供了真实材料在 HUVEC 中的响应曲线，展示了时间变化与条件相关的材料差异。两部分共同说明，预测性能和实验响应结构需要各自获得直接证据，方向一致的实验趋势不足以替代逐样本预测评价。

% Word body block 100
化学增强的阴性结果限定于当前字段质量、样本组成与模型配置，并不说明材料化学性质在毒性形成中没有作用。历史数据的 Zeta 电位缺失率为 42.40\%，修饰物类别和分子量没有完整字段，晶体信息编码又只反映是否报告晶体学信息，这些条件限制了模型对材料差异的识别。表征方法、报告习惯和研究环境还可能共同影响输入值，新增字段是否具有生物学意义，与其能否在当前归档数据中提供稳定预测增量，属于需要分开讨论的问题。

% Word body block 101
同样，迁移对照尚未建立稳定收益，可能与不同材料体系的响应分布、共同特征的覆盖以及当前树模型迁移接口有关。RF 的新增树训练保留原有通用树，XGBoost 的续训依赖已学习的通用结构，本研究没有检验其他迁移方法，不能将结果推广为迁移学习对 SeNPs 普遍无效。对于现阶段的材料筛选，更合适的工作是明确模型可接受的输入、记录外部测试条件，并检验预测误差，而不是根据通用基准指标直接判断新材料预测精度。

\subsection{与既有计算研究的关系}

% Word body block 103
Gul 等的跨材料属性分析与 Masarkar 等的预测研究，为组织纳米材料理化性质、细胞背景和暴露条件提供了基础 [3, 4]。本研究进一步关注这些信息在 SeNPs 文献分组评价下的增量，以及新增实验如何与历史建模衔接。由于数据纳入、特征定义、模型选择与验证划分均不同，本研究不以跨论文的单个 R2 或 MAE 数值宣称性能优于既有工作，也不将归档中缺少完整训练证据的深度模型指标纳入当前主比较。

% Word body block 104
嵌套评价使内层选择与外层报告分离，对有限样本下模型选择偏差的控制与已有方法学研究一致 [5]。不过，文献分组验证只检验对未参与训练的文献组的表现，不能自动代表对未知配方、未知修饰物或不同实验室的泛化。旧历史复现还混合了随机交叉验证误差、全量拟合 R2 和不同预处理设置，当前结果不构成只改变划分方式的受控对照，因此性能差异不能全部归因于随机划分造成的泄漏，方法与数据口径的变化也需要保留在解释中。

\subsection{实验观测的意义与机制解释范围}

% Word body block 106
HUVEC 数据把研究从历史文献预测扩展到真实材料的剂量—时间观测，最高测试浓度下的时间变化与有限 IC50 估计数量的增加共同表明，评价时间会影响材料差异的呈现方式。对照已有SeNPs 细胞研究，这些结果支持将材料修饰和暴露条件一起纳入实验比较 [1, 2]，但当前单一细胞模型不能提供正常与肿瘤细胞选择性、跨细胞普适性或体内安全性的证据。

% Word body block 107
材料的高浓度排序与表观 IC50 排序并不一致，反映了固定浓度残余存活率与跨越 50\% 阈值所处位置的不同。表观 IC50 仅由两个相邻报告点插值，浓度间距、局部波动和检测读数都可能影响估计，尚不能用于给出精细的材料优劣判断。对没有跨越阈值的曲线保留 > 512 状态，也避免了在测试范围外强行拟合一个看似完整的数值。

% Word body block 108
粒径、电位和修饰信息在不同材料之间同时变化，每种材料又仅有一组描述符报告值，因此弱秩相关既不能证明这些性质无作用，也不能将某种材料的低存活率归因于单个描述符。大于 100\%的报告值同样需要谨慎解释，当前缺少原始孔读数、空白扣除及检测干扰信息，不能据此判定真实增殖促进或特定毒性机制。核对这些实验信息后，才能进一步评估剂量曲线的波动与材料相关的检测影响。

\subsection{外部预测评价仍需完成的环节}

% Word body block 110
冻结模型的逐条件预测文件具有解盲前封存记录，封存时间为 2026 年 10 月 13 日 16:23:13，记录的 SHA-256 与当前 model\_predictions.csv 一致。按材料、报告浓度、时间及浓度单位匹配后获得 351 对预测—实测值，因而本研究可以报告当前条件表上的 MAE、RMSE、平均偏差及预测—实测排序一致性。该评价仍受浓度单位及换算基准是否正确的影响，且每个条件只有一个归档报告值，不能据此估计独立实验重复间的变异。实测 IC50 不能作为新材料初筛模型的输入。
